Praktický krok‑za‑krok plán pro zavedení AI do kurzu, s důrazem na řízení rizik a kontrolu kvality.

\begin{enumerate}
\item Identifikujte cíle a scénáře (kde AI přidává hodnotu, měřitelné metriky).
\item Návrh malého pilotu: cíl, rozsah, KPI, role a komunikace se studenty.
\item Kontrola kvality a audit: rubriky (A–D, viz 3.3); auditní záznamy (ID, datum, verze rubriky, prompt/LLM skóre, lidské skóre, důvod eskalace, rozhodnutí).
\item Řízení rizik: fact‑check vrstvy, detekce plagiátorství, konzervativní prompty, minimalizace sdílených dat a zajištění souhlasu.
\item Školení a zapojení studentů: krátké školení pro vyučující/hodnotitele; transparentnost a sběr zpětné vazby.
\item Monitorování a iterace: KPI (shoda AI vs. člověk, čas na feedback, spokojenost, eskalace); pravidelné revize a úpravy workflowů.
\end{enumerate}

Krátká šablona pilotního plánu:
\begin{description}
\item[Cíl pilotu:] $\underline{\hspace{6cm}}$ \item[Rozsah:] $\underline{\hspace{6cm}}$ \item[KPI:] $\underline{\hspace{3cm}}$ \item[Role:] $\underline{\hspace{3cm}}$
\item[Kontrola kvality / auditní záznam:] (povinné položky viz výše) \item[Plán komunikace \\ \& Harmonogram:] $\underline{\hspace{6cm}}$
\end{description}
Poznámka: pro konkrétní rubriky a auditní záznamy viz sekce 3.3.